\documentclass[a4paper,12pt]{article}
\usepackage[utf8]{inputenc}
\usepackage[ngerman]{babel}
\usepackage{amsmath}
\usepackage{geometry}
\geometry{a4paper, top=15mm, left=15mm, right=15mm, bottom=15mm,
	headsep=10mm, footskip=12mm}
\usepackage{parskip}
\usepackage{circuitikz}
\usepackage{graphicx}
\usepackage{tikz}
\usepackage{href-ul}
\hypersetup{
	colorlinks=true,
	linkcolor=blue,
	urlcolor=blue}

\begin{document}
	
	\begin{tikzpicture}(10,3)
		\draw[ultra thick](2,0) --(10,0) -- (10,1.3) --(2,1.3) -- (2,0);
		\draw[fill=black](2,0)-- (0,.6) -- (2,.6) -- (2,0);
		\node at (6,.6) {\large Lösungsblatt von \href{https://www.okuyakl.de}{www.okuyakl.de}};
	\end{tikzpicture}
	
	\section*{Grundwissen – Lösungen}
	
	\begin{itemize}
		\item \textbf{Ein geschlossener Stromkreis} ist notwendig, damit eine Lampe \textbf{leuchtet}.
		
		\item Der Strom fließt vom \textbf{Plus}-Pol der Batterie zum \textbf{Minus}-Pol (technische Stromrichtung).
		
		\item In einem Stromkreis braucht man mindestens:
		\begin{itemize}
			\item( x ) eine Batterie oder Spannungsquelle
			\item( x ) einen Verbraucher (z. B. Lampe)
			\item( x ) Kabel oder Leitungen
		\end{itemize}
		
		\item Symbole:
		
		\textbf{Batterie:}
		\begin{circuitikz}
			\draw (0,0) to[battery] (2,0);
		\end{circuitikz}
		
		\textbf{Lampe:}
		\begin{circuitikz}
			\draw (0,0) to[lamp] (2,0);
		\end{circuitikz}
		
		\textbf{Schalter (offen):}
		\begin{circuitikz}
		 \draw (0,0) to[switch] (2,0);
		\end{circuitikz}
	\end{itemize}
	
	\hrule
	
	\section*{1. Stromkreis verstehen }
	

\begin{itemize}
	\item a) Warum muss der Stromkreis „geschlossen“ sein? - Damit die Elektronen ohne Hindernis fließen können. Die Spannungsquelle arbeitet wie eine Elektronenpumpe und erzeugt am Pluspol einen Elektronenmangel und am Minuspol einen Elektronenüberschuss.
	
	\vspace{1em}
	\underline{\hspace{12cm}}
	
	\item b) Welche Möglichkeiten gibt es für die  Schaltung? 
	
	
\begin{minipage}{0.4 \pagetotal}
	
	 Reihenschaltung:
		\begin{circuitikz}[american voltages]
			% Batterie links, zwei Lampen oben in Reihe, Rückleitung unten
			\draw
			(0,0) to[battery1,l=$V$] (0,3)   % Spannungsquelle
			-- (1.5,3) to[lamp,l=$L_1$] (2.5,3) % erste Glühbirne
			to[lamp,l=$L_2$] (6,3)            % zweite Glühbirne
			-- (6,0) -- (0,0);                % Rückleitung zum Minuspol
		\end{circuitikz}
		
\end{minipage}
\begin{minipage}{0.4 \pagetotal}
	
	Parallelschaltung:
	\begin{circuitikz}[american voltages]
		% Batterie links
		\draw
		(0,0) to[battery1,l=$V$] (0,4)
		-- (4,4)                % obere Leitung
		(0,0) -- (4,0);         % untere Leitung
		
		% Zwei Lampen parallel (senkrecht)
		\draw
		(2,4) to[lamp,l=$L_1$] (2,0);
		\draw
		(4,4) to[lamp,l=$L_2$] (4,0);
	\end{circuitikz}
\end{minipage}

\newpage
	\item c) Kreuze an:  
	In einer {\bf Reihenschaltung} zweier Lampen...
	
	\begin{itemize}
		\item( ) kann man beide Lampen unabhängig voneinander aus- und einschalten.
		\item(x) gehen beide Lampen aus, wenn eine durchbrennt.
		\item(x) fließt derselbe Strom durch alle Lampen, auch wenn sie verschieden sind.
		\item( ) wird es heller, wenn man mehr Lampen einbaut.
		\item(x) müssen beide Lampen die gleiche Betriebsstromstärke in Ampere haben.
		\item( ) müssen beide Lampen die gleiche Betriebsspannung in Volt haben.
		\item( ) müssen beide Lampen die gleiche Leistung  in Watt haben.
		\item(x) muss die Summe der Betriebsspannungen der Lampen gleich der Spannung der Stromquelle sein.
	\end{itemize}
	
	\item d) Kreuze an:  
	In einer {\bf Parallelschaltung} zweier Lampen...
	
	\begin{itemize}
		\item(x) kann man beide Lampen unabhängig voneinander aus- und einschalten.
		\item( ) gehen beide Lampen aus, wenn eine durchbrennt.
		\item( ) fließt derselbe Strom durch alle Lampen, auch wenn sie verschieden sind.
		\item(x) wird es heller, wenn man mehr Lampen einbaut.
		\item( ) müssen beide Lampen die gleiche Betriebsstromstärke in Ampere haben.
		\item(x) müssen beide Lampen die gleiche Betriebsspannung in Volt haben.
		\item( ) müssen beide Lampen die gleiche Leistung  in Watt haben.
		\item( ) muss die Summe der Betriebsspannungen der Lampen gleich der Spannung der Stromquelle sein.
	\end{itemize}
	
\begin{minipage}{0.45\textwidth}
	\item e) {\bf Reihenschaltung} Der Stromkreis ist unverzweigt. Wenn ein Schalter in der Reihe geschlossen wird, gehen beide Lampen an. Und wenn er geöffnet wird, gehen beide Lampen aus.
\end{minipage}
\hspace{1 cm}
\begin{minipage}{0.45\textwidth}
 
	\begin{circuitikz}[american voltages]
		% Batterie links, zwei Lampen oben in Reihe, Rückleitung unten mit Schalter
		\draw
		(0,0) to[battery1,l=$V$] (0,3)   % Spannungsquelle
		-- (1.5,3) to[lamp,l=$L_1$] (2.5,3) % erste Glühbirne
		to[lamp,l=$L_2$] (6,3)            % zweite Glühbirne
		-- (6,0)
		to[switch,l=$S$] (1,0)            % Schalter
		-- (0,0);                         % Rückleitung zum Minuspol
	\end{circuitikz}
\end{minipage}
	
\begin{minipage}{0.45\textwidth}
	{\bf Parallelschaltung} Der Stromkreis teilt sich in zwei Äste. Zu jeder Lampe kann ein Schalter eingebaut werden. Mit Schalter 1 kann Lampe 1 an- und ausgeschaltet werden, und mit Schalter 2 unabhängig davon Lampe 2.
\end{minipage}
\hspace{1 cm}
\begin{minipage}{0.45\textwidth}
			\begin{circuitikz}[american voltages]
				% Batterie links, obere und untere Sammelschiene
				\draw
				(0,0) to[battery1,l=$V$] (0,4)
				-- (3.5,4)
				(0,0) -- (3.5,0);
				
				% Linker Zweig: Schalter + Lampe L1
				\draw
				(1.5,4) to[lamp,l=$L_1$] (1.5,2)
				to[switch,l=$S_1$] (1.5,0);
				
				% Rechter Zweig: Schalter + Lampe L2
				\draw
				(3.5,4) to[lamp,l=$L_2$] (3.5,2)
				to[switch,l=$S_2$] (3.5,0);
			\end{circuitikz}
	\end{minipage}
\end{itemize}

\begin{center}
	\includegraphics[width=7 cm]{../../viecher/endcomic.pdf}
	
	Hier geht es zurück zum \href{https://www.okuyakl.de/phys/p7escL129/aa129.pdf}{Aufgabenblatt}
\end{center}
\end{document}
